% Full mathematical argument for the separately frozen Phase B algorithm.
% This source was prepared without scientific execution or typesetting.
\section{Qualification and work of the bounded preliminary service}
\label{supp:budgetedtail119}

\subsection{Objects, information, and algorithm}
Fix a state $x$, its original feasible action interval $[0,A(x)]$, a
realized proposal $a\in[0,A(x)]$, and
$\alpha\in\{1/2,3/4\}$.  Write
\[
\begin{gathered}
 R=1/32,\qquad t=1-\alpha,\qquad L=2R\alpha,\qquad I=[-R,R-L],\\
 f_a(z)=V_1(F(x,a,z)).
\end{gathered}
\]
The primitive transition is affine in the innovation and the action.
The already established original bounds give convexity of $f_a$,
its strong-convexity modulus $m_z=181/64$, and
$V_1\le J_1\le V_1+\varepsilon_1$ for the implemented terminal
controller, with $\varepsilon_1=2^{-20}$.  Its complete-policy cost is
\[
 J(a)=S(x)+h(a)+\beta\operatorname{AVaR}_\alpha(J_1\circ F_a),
\]
whereas $Q(a)$ uses $V_1$ in this expression, and
$V_0(x)=\inf_{0\le v\le A(x)}Q(v)$.

Two information types are distinguished.  A floating terminal value
is used only to propose a cut.  An outward terminal payload encloses
the original $V_1$ value; when used at a quadrature midpoint it also
supplies a valid global affine minorant, with every uncertain slope
component retained.  These are the same original payloads used by the
common terminal compiler.  Counting a payload does not assign unit
cost to its internal original-risk evaluation, optimization, interval
arithmetic, or supporting-witness construction.

The floating proposal compares $f_a(u+L)-f_a(u)$ at the two endpoints
of $I$.  A proposed endpoint exit suppresses further floating search
for that row.  Otherwise six interval-halving updates are performed,
after which the midpoint of the final proposed bracket is rounded and
clipped to the feasible $2^{-32}$ grid.  The signs need not be correct:
only feasibility of the final cut $u$ is used below.  The production
domain checks enforce this feasibility and the declared risk levels.

At this $u$, two outward payloads give an interval $[d_-,d_+]$
containing $f_a(u+L)-f_a(u)$.  The computed outward upper bound is
\begin{equation}\label{eq:budgetedkappa119}
 \widehat\kappa\ge
 \frac{1}{2Rt}\max_{d\in\{d_-,d_+\}}
 \left[-\min_{v\in I}
       \left\{d(v-u)+\frac{m_zL}{2}(v-u)^2\right\}\right].
\end{equation}
The same state, action, risk level, cut, and defect are cached for the
two attempted resolutions.  Reusing the cache for another state,
action, or objective is disallowed.

Let the lengths of the left and right tails be
$b_1=u+R$ and $b_2=R-u-L$.  For $m=2$ and then, if needed, $m=4$,
partition each tail with $b_j>0$ into $m$ equal cells.  Each cell of
that tail has weight $b_j/(m\,2Rt)$.  A cell uses one midpoint
payload and its two endpoint upper values; endpoints shared by
neighboring cells within the same attempted resolution are acquired
once.  No values are requested solely for a zero-length tail.  At
each resolution form the original-policy and original-optimum bounds
of Theorem~\ref{thm:investmenttail119}; accept only if the outward
regret allowance is at most $1/256$ and the outward policy width is
at most $1/8192$.  If neither attempt qualifies, continue with the
frozen 128-bin service from the already acquired proposal.  Its
original conventional repair, further resolutions, and rational
fallback are included.

\subsection{Proof of Proposition~\ref{prop:budgetedtail119}}
\begin{proof}
We establish feasibility, validity of each possible acceptance, and
the full preliminary query count separately.

\paragraph{Exact zero-mass omission.}
The tail lengths are nonnegative and satisfy
\[
 b_1+b_2=(u+R)+(R-u-L)=2R-L=2Rt>0.
\]
Consequently the number $p$ of positive-length tails is either one
or two, and the retained weights sum to one.  The density relative
to the original uniform law is still $t^{-1}$ on these tails and
zero on the omitted central interval.  Removing a zero-length tail
removes a set of zero original probability; it changes none of the
exact integrals, risk witnesses, or action supports.  This is not an
approximation by a small-mass cutoff.

In the implemented domain $u$ is a feasible multiple of $2^{-32}$,
$R=2^{-5}$, and $L$ and $2Rt$ are dyadic.  With dyadic $m\le4096$,
the tail endpoints, midpoints, and weights used here have exact
binary representations with fewer than 53 significant bits.  To see
this directly, write $u=k2^{-32}$ and $m=2^q$, with $q\le12$.
The nonnegative tail lengths have integer numerators at most $2^{27}$
on that grid.  Multiplying them by a grid fraction $j2^{-q}$ gives
numerators at most $2^{39}$ on the $2^{-(32+q)}$ grid.  Adding a
tail endpoint and then taking a midpoint needs at most 41 significant
binary digits.  Dividing a tail length by $m\,2Rt$ only shifts its
binary exponent, because $2Rt$ is either $2^{-5}$ or $2^{-6}$.
All these values remain in the normal binary64 exponent range.  Thus
the production predicate that a weight is positive is precisely a
positive-length predicate.  The three groups of rows with only a
left tail, only a right tail, or both tails partition the admissible
batch.  Within each group, the retained endpoint columns belong to
the same successive cells as the retained midpoint and weight
columns.  Grouping therefore leaves the preceding exact weighted
objects unchanged.  Floating summation orders can alter outward
interval widths, but must still enclose those objects; bitwise
equality of separately implemented certificates is not assumed.

\paragraph{The proposed cut remains certified by the original continuation.}
Let $H_a(v)=\int_v^{v+L}f_a(z)\,dz$.  The original curvature bound
implies that $H_a$ is $m_zL$-strongly convex.  For the true
$d=f_a(u+L)-f_a(u)$ and every feasible $v$,
\[
 H_a(v)\ge H_a(u)+d(v-u)+\frac{m_zL}{2}(v-u)^2.
\]
Taking the infimum over the entire original interval $I$ bounds
$H_a(u)-\min_IH_a$.  The negative of the quadratic minimum in
\eqref{eq:budgetedkappa119} is convex as a function of $d$.
Its value at any $d\in[d_-,d_+]$ is therefore no larger than the
larger endpoint value.  An outward lower bound on the constrained
quadratic minimum, followed by the indicated upper operations,
preserves the inequality.  Finally,
\[
 \operatorname{AVaR}_\alpha(f_a)-R_u(f_a)
       =\frac{H_a(u)-\min_I H_a}{2Rt}
       \le\widehat\kappa.
\]
This derivation uses no floating sign and no assertion that the cut
optimizes the original risk.  It applies to either endpoint, an
incorrect floating endpoint exit, and any imperfect interior cut.

\paragraph{Original policy and all-action validity.}
For a positive-length cell $B_i=[l_i,r_i]$ with midpoint $z_i$,
write $y_i=F(x,a,z_i)$ and let its valid global affine witness be
\[
 V_1(y)\ge\ell_i+p_i^{\mathsf T}(y-y_i).
\]
Integrating over the cell cancels the innovation term proportional
to $z-z_i$.  If $w_i=(r_i-l_i)/(2Rt)$, this proves
$r_-:=\sum_iw_i\ell_i\le R_u(f_a)$.  Convexity of $f_a$ and
the outward endpoint upper values prove
$R_u(f_a)\le r_+:=\sum_iw_i(u_i^-+u_i^+)/2$.
The cut-defect result and the unchanged terminal controller then give
\[
 S(x)+h(a)+\beta r_-
 \le J(a)
 \le S(x)+h(a)+\beta\{r_++\widehat\kappa+\varepsilon_1\}.
\]
The last term follows from monotonicity and cash invariance applied
to $V_1\le J_1\le V_1+\varepsilon_1$.  No convexity of the
implemented $J_1$ is needed.

For any other original feasible action $v$, use the same fixed
density and the same global supports.  The action-exposure vector
$b$ satisfies $F(x,v,z)=F(x,a,z)+b(v-a)$.  Integration and
$h''\ge2$ yield
\[
\begin{gathered}
 Q(v)\ge W+s(v-a)+(v-a)^2,\qquad W=S(x)+h(a)+\beta r_-,\\
 s=h'(a)+\beta\sum_iw_i b^{\mathsf T}p_i.
\end{gathered}
\]
Taking the infimum over $[0,A(x)]$ gives a lower bound on
$V_0(x)$.  If $s$ is interval-valued, the constrained minimum is
concave in $s$, so the smaller minimum at its two enclosing
endpoints remains a valid lower bound.  This is the original
all-action construction, with every arithmetic endpoint directed
outward.  The policy upper bound minus this optimum lower bound
therefore encloses complete-policy regret, and the two policy
endpoints enclose policy cost.  The qualification test at either
attempt preserves precisely the two stated original targets.
No monotonicity of successive reported widths is required.

If both tests fail, the preliminary certificates do not overwrite
the actor proposal, the economic primitives, the deployed future
policy, or the frozen fallback.  The continuation branch first
constructs its 128-bin original certificate and then runs the
unchanged repair and refinement sequence.  Its own outward
qualification is required before return.  A finite prelude of
terminating oracle calls therefore preserves the termination and
qualification guarantee of the underlying complete service.  A
computation that fails an enclosure or qualification check must
retain its failure evidence and must not return a successful
scientific decision record.

\paragraph{All preliminary values are counted.}
The two floating endpoint differences request four terminal values
per proposed action, including any repeated locations.  Only rows
without a proposed endpoint exit receive further queries.  Six
such updates request two values each, so their total is at most
$4+6\cdot2=16$; a proposed endpoint exit uses four.

The outward defect requests exactly two values once for that
state/action/risk identity.  A resolution with $m$ cells in each
of $p$ nonempty tails uses $pm$ midpoint payloads and $p(m+1)$
endpoint payloads, hence $p(2m+1)$.  The calculation allows repeated
locations across resolutions, between a midpoint and an earlier
endpoint, and between the defect and quadrature.  All such repeated
calls remain charged.  A first-resolution acceptance thus uses at
most $2+5p$ outward payloads.  If the second resolution is attempted,
the complete preliminary total is
\[
 2+p\{(2\cdot2+1)+(2\cdot4+1)\}
      =2+14p\le30.
\]
The cache is why the two defect values are not reacquired at the
second resolution.  Grouping rows according to their positive tails
changes the number of vector calls, not these scalar-payload counts.
An earlier failed call cannot enlarge the prescribed subsequent
schedule; it terminates with a failed record instead of licensing
an unrecorded retry.

Every call made by the frozen continuation is additional to this
preliminary count.  If its realized continuation transcript has
$N_F^{\rm outward}$ outward and $N_F^{\rm floating}$ floating
terminal payloads, the corresponding combined counts on a branch
that invokes it are bounded by
$30+N_F^{\rm outward}$ and $16+N_F^{\rm floating}$.
Actor acquisition and all non-payload operations are outside these
two counts and remain in the complete-work ledger.  This proves
the proposition without assuming a wall-clock cost per payload or
an economic ordering of the returned actions.
\end{proof}

\subsection{Ablation and prospective scope}
The masked-tail ablation uses the same proposal, cache, exact
zero-mass omission, original bounds, and fallback, while attempting
all $m\in\{2,4,8,16,32,64\}$.  The preceding proof applies to each
individual acceptance.  Its full preliminary outward bound is
\[
 2+p\sum_{m\in\{2,4,8,16,32,64\}}(2m+1)
      =2+258p\le518,
\]
with the same floating bound of 16.  Its additional attempted
resolutions are not free, even if the later fallback is identical.

These two schedules were chosen after the separately preserved
Phase A results and are therefore declared as Phase B algorithm
development.  The five-service comparison reacquires the common
training information and executes its conventional controls and
original-policy reference enclosures afresh under its own source
freeze.  Neither the mathematical validity of early acceptance nor
the information cap establishes lower elapsed time, lower policy
regret, a calibrated adoption decision, or neural specificity.
Those claims require the corresponding current records.  In
particular, early acceptance may keep an actor proposal that another
service repairs; its complete original-policy cost must be evaluated
as that returned policy, with the unchanged future controller.
